\section{Methodology}
\label{sec:methodology}

To measure whether static and execution feedback detect orthogonal error classes, we design an analysis that extracts error sets from both sources on the same code samples and measures their overlap.

\subsection{Overview}

Our approach has three components: (1) static analysis pipeline to extract structural errors, (2) execution pipeline to extract behavioral errors, and (3) orthogonality measurement using Jaccard similarity.

\subsection{Static Analysis Pipeline}

We use pylint and mypy, standard Python linters that together cover type errors, syntax errors, undefined variables, unused imports, and other structural issues.

\textbf{Tool Configuration:}
\begin{itemize}
    \item pylint: Error (E) and Warning (W) codes
    \item mypy: \texttt{--ignore-missing-imports} flag
\end{itemize}

For each code sample, we extract the set of error categories detected, not individual instances.

\subsection{Execution Pipeline}

We use the EvalPlus test harness with extended test suites. HumanEval+ provides 80x more tests than original; MBPP+ provides 35x more.

\textbf{Configuration:} Timeout 5s per problem, 8 parallel workers.

\subsection{Orthogonality Measurement}

We measure overlap using Jaccard similarity:
\begin{equation}
J(S, E) = \frac{|S \cap E|}{|S \cup E|}
\end{equation}
where $S$ = problems with static errors, $E$ = problems with execution errors. $J = 0$ indicates complete orthogonality.

\subsection{Analysis Framework}

We decompose into three sub-hypotheses:
\begin{itemize}
    \item \textbf{H-E1:} Jaccard $< 0.3$ (orthogonality)
    \item \textbf{H-M1:} Static coverage $> 60\%$
    \item \textbf{H-M2:} Behavioral rate $> 40\%$
\end{itemize}
